\chapter{Emergence and Framework Formation}
\label{ch:emergence}

\section{Why This Chapter Must Become Stricter}

The earlier version of this chapter correctly saw that frameworks are formed,
not merely encountered. But it still read too much like a mechanism sketch. If
the manuscript wants to claim that habits, interpretive frames, and behavioural
regimes are historically produced structures, then it must explain
\emph{how} higher-level organization arises, \emph{what} kind of evidence
supports that claim, and \emph{where} the analogy to human behaviour becomes
interpretive rather than directly measured.

This chapter therefore adopts a visible three-layer architecture:
\begin{enumerate}[nosep]
  \item \textbf{Biological grounding:} what current biology and neuroscience
        allow us to say about organized state formation and stability.
  \item \textbf{Physical and systems analogies:} what complex-systems and
        self-organization language contributes when used carefully.
  \item \textbf{Mathematical and explanatory structure:} what emergence,
        path dependence, metastability, and multiscale organization mean in a
        disciplined model.
\end{enumerate}

The evidence policy remains explicit. The chapter's main direct sources are
\emph{Metastability demystified: the foundational past, the pragmatic present
and the promising future} (Nature Reviews Neuroscience, 2024, \textbf{medium}),
\emph{Associative conditioning in gene regulatory network models increases
integrative causal emergence} (Communications Biology, 2025, \textbf{medium}),
and the neighbouring empirical consciousness chapter, whose strong sources
justify taking state structure, stability, and transition seriously at the
biological level. The chapter also leans on classic textbooks such as John
Holland's \emph{Emergence}, \emph{Introduction to the Theory of Complex
Systems}, and mathematically oriented dynamical-systems texts. None of these
sources directly proves the full behavioural framework of this book. Their
value is to constrain the explanation so it is structurally serious rather than
metaphorical improvisation.

\section{What Emergence Means, and What It Does Not Mean}

\begin{definition}[Emergence]
\textbf{Emergence} is the appearance of a higher-level organized pattern whose
behaviour cannot be usefully described as a mere list of isolated local events,
because the pattern feeds back on later local behaviour, alters future costs,
or becomes an explanatory layer in its own right.
\end{definition}

\begin{definition}[Framework formation]
\textbf{Framework formation} is the process by which repeated inputs,
conditioning, environmental coupling, and historical reuse generate a
relatively stable behavioural or cognitive regime.
\end{definition}

\begin{remark}
The chapter uses ``emergence'' in a disciplined explanatory sense, not as a
magic word for anything complicated. A pattern is not emergent merely because
it is surprising, difficult, or emotionally important.
\end{remark}

\subsection{What Emergence Does Not Mean}

For this manuscript, emergence does \emph{not} mean:
\begin{enumerate}[nosep]
  \item \textbf{Mystery inflation:} calling a process emergent because its
        micro-details are inconvenient to track.
  \item \textbf{Causal vagueness:} pretending that once a macro-pattern
        appears, local mechanisms no longer matter.
  \item \textbf{Unrestricted analogy:} assuming that because emergence occurs
        in one biological or physical system, the same mechanism must govern a
        particular human routine one-to-one.
  \item \textbf{Strong proof by vocabulary:} upgrading system language such as
        ``basin,'' ``phase,'' or ``conditioning'' into direct evidence about
        daily agency without an empirical bridge.
\end{enumerate}

\begin{discussion}
This negative definition is strategically important. Without it, the chapter
would invite exactly the criticism it should avoid: that emergence is being
used as a prestige label rather than as a constrained explanatory tool. The
book's claim is narrower. Repeated local interaction can generate
history-sensitive organization; that organization can later structure what is
easy, salient, and likely; and this makes the higher-level regime worth naming.
\end{discussion}

\section{Layer I: Biological Grounding}

\subsection{Why Biology Belongs In A Chapter On Formation}

The neighbouring consciousness chapter established one strong empirical floor:
biological systems exhibit organized state signatures, stability structure, and
nontrivial transition dynamics. That result does not yet prove ``framework
formation'' in the everyday sense, but it changes what is intellectually
available. The book is now allowed to ask not only how states persist, but how
repeated coupling and conditioning build regimes that later alter transition
cost.

\begin{proposition}[History-sensitive inertia]
The inertia of a current state depends not only on the present stimulus, but
also on the accumulated reinforcement history, expected reward landscape, and
prior successful re-entry routes associated with that state.
\end{proposition}

\begin{discussion}
This proposition is not meant as a direct theorem extracted from one paper. It
is a synthesis claim justified at \textbf{medium} strength by the chapter's
sources. The empirical consciousness papers support taking structured states
seriously. The causal-emergence and metastability sources support treating
history and system organization as explanatory variables rather than merely
decorative background.
\end{discussion}

\subsection{Conditioning and Integrative Organization}

The 2025 Communications Biology paper on associative conditioning in gene
regulatory network models is especially useful because it gives a tractable
example of repeated local interaction producing stronger higher-level
organization. This source should be read carefully. It does \emph{not} show
that human study routines are literally gene regulatory networks. What it buys
the manuscript is more limited and more valuable:
\begin{enumerate}[nosep]
  \item repeated local couplings can alter global organization,
  \item higher-level organization can acquire explanatory value,
  \item conditioning history can change what later transitions look like.
\end{enumerate}

For framework formation, this means a regime is not merely the sum of discrete
past acts. Once cues, rewards, expectations, and environmental reinforcements
become mutually stabilizing, the resulting pattern becomes a real operating
layer. In business terms, the system begins to ``carry policy'' in its own
structure.

\subsection{Metastability Rather Than Simple Rigidity}

The 2024 Nature Reviews Neuroscience article on metastability is critical
because it blocks a false binary. Frameworks are often described as either
rigid habits or flexible choices. Metastability offers a better middle term:
organized patterns can be durable enough to persist and yet reconfigurable
enough to release under changed conditions.

\begin{definition}[Metastable framework]
A \textbf{metastable framework} is a regime that is locally persistent under
ordinary perturbations, but not globally fixed against larger contextual,
reward, or control-parameter changes.
\end{definition}

\begin{discussion}
This definition matters because many productive frameworks must be metastable,
not absolute. A study routine that can never be interrupted is dysfunctional.
A routine that collapses under every disturbance is useless. The chapter
therefore treats metastability as the right biological and systems-level bridge
for explaining why good frameworks feel stable without becoming cages.
\end{discussion}

\section{Layer II: Physical and Systems Analogies}

\subsection{Why Analogy Is Needed, But Must Stay Disciplined}

Biology gives the chapter legitimacy, but not all the conceptual machinery.
Complex-systems and physical analogies are useful because they provide
organized ways to think about repeated interaction, cooperative locking,
threshold effects, and multiscale order. Their role in this chapter is
\textbf{medium} and explicitly interpretive.

\begin{remark}
The manuscript does not claim that a study desk, a research lab, and a neural
assembly are the same object. It claims that they may share analyzable
properties such as coupling, reinforcement, path dependence, and reorganized
cost landscapes.
\end{remark}

\subsection{Self-Organization and Order Parameters}

Classic complexity texts are useful here because they explain a foundational
idea: local interaction can produce higher-level coordination without requiring
a central executive that explicitly computes the full pattern in advance.

In the language of this manuscript, a framework can be treated as an
order-parameter-like summary of many local regularities:
\begin{itemize}[nosep]
  \item time of day,
  \item cue placement,
  \item posture,
  \item reward expectation,
  \item social synchronization,
  \item memory of prior success and prior failure.
\end{itemize}

Once these variables align often enough, the system no longer has to negotiate
the full behavioural choice from scratch. A macro-pattern exists, and that
pattern constrains the next round of micro-decisions.

\subsection{Path Dependence and Interference}

Path dependence is essential to the user requirement about ``past data
interference.'' The same present stimulus can lead to different action because
the system is carrying different traces of prior conditioning, prior losses,
prior rewards, or prior environmental pairings.

\begin{definition}[Path dependence]
\textbf{Path dependence} means that present transition costs depend on the
sequence of prior states and reinforcements, not only on the current snapshot.
\end{definition}

\begin{definition}[Interference]
\textbf{Interference} is the distortion or blocking of a target transition by
traces of previously reinforced routes, cues, or expectations that remain
active in the current context.
\end{definition}

\begin{discussion}
This is one of the chapter's most business-relevant points. Bad frameworks are
often not recreated because the present decision is morally wrong; they are
recreated because the old route still has more stored structure. Likewise, good
frameworks become easier not because the agent becomes ``better'' in the
abstract, but because the path has been rendered cheaper through repeated use.
\end{discussion}

\subsection{Multiscale Organization}

Emergence also matters because organization exists at several scales at once.
A useful framework can be:
\begin{enumerate}[nosep]
  \item locally enacted by small repeated actions,
  \item maintained by room layout, schedule structure, or social timing,
  \item embedded inside larger biological and conscious-state conditions.
\end{enumerate}

This multiscale viewpoint prevents a common mistake: trying to explain all
framework formation either by inner psychology alone or by environment alone.
The book's claim is more demanding. Frameworks form when several scales align
long enough to stabilize one route over others.

\section{Layer III: Mathematical and Explanatory Structure}

\subsection{From Local Updates To Global Regimes}

The mathematical value of emergence language is not that it creates exact
closed-form formulas for everyday life. Its value is that it clarifies the
architecture of explanation.

\begin{proposition}[Micro-to-macro feedback]
If repeated local updates modify the cost of future transitions, then the
resulting macro-pattern is not merely descriptive. It becomes causally
relevant within the model because it changes what later trajectories are cheap,
stable, or accessible.
\end{proposition}

\begin{discussion}
This is the explanatory point behind causal-emergence language. Once a mature
regime changes later transition geometry, the regime deserves to be treated as
more than a summary label. It has become a control-relevant level of
description.
\end{discussion}

\subsection{Regimes, Basins, and Re-entry Cost}

The chapter does not require heavy equations, but it does require a stricter
mathematical picture. Let a behavioural regime be represented as a region of
state space with the following properties:
\begin{enumerate}[nosep]
  \item \textbf{entry cost:} how much coordination is required to enter it;
  \item \textbf{retention:} how likely trajectories are to remain inside it
        once entered;
  \item \textbf{re-entry economy:} how cheaply the system can return after
        interruption;
  \item \textbf{competitive suppression:} how strongly rival routes are made
        less attractive by the established regime.
\end{enumerate}

Framework formation can then be interpreted as a process that progressively:
\begin{itemize}[nosep]
  \item lowers entry cost for one route,
  \item deepens or sharpens its basin of attraction,
  \item stores route memory through repeated coordinated use,
  \item increases the interference cost of competing paths.
\end{itemize}

\subsection{Why Emergence Is Not Just Aggregation}

Simple aggregation would mean that many repetitions merely add up. Emergence in
the present chapter is stronger than that. The repeated acts do not just
accumulate; they reorganize the future problem.

\begin{lemma}[Reorganization criterion]
If repeated local events change the future transition landscape itself, then
the resulting higher-level pattern is doing more than summarizing the events.
\end{lemma}

\begin{proof}
Suppose the local events merely add up without reorganizing later choice. Then
the future decision problem is structurally unchanged, and the repeated events
have not created a new explanatory level. But if the events lower one route's
re-entry cost, strengthen its retention, or increase interference against
rival routes, then the future problem is different from what it would have been
without the history. The macro-pattern has therefore become model-relevant.
\end{proof}

\subsection{Mathematical Upgrade Path}

This chapter is intentionally lighter than Chapter~\ref{ch:advanced-dynamics},
but it prepares that chapter in three ways:
\begin{enumerate}[nosep]
  \item metastability clarifies why useful frameworks must persist without
        becoming immutable;
  \item attractor language clarifies how repeated use can turn one route into a
        default basin;
  \item operator and control viewpoints later help describe how interventions
        act on a history-sensitive system rather than a blank slate.
\end{enumerate}

\section{Example Bank: How Frameworks Actually Form}

\subsection{Example 1: Study Entry Ritual}

\begin{example}[Repeated entry cue creates cheap re-entry]
A student repeatedly begins work at the same place, with the same notebook
opened to the next unresolved line, and the same first five-minute task. Over
time the study state becomes easier to enter after small interruptions.
\end{example}

\begin{discussion}
This is a clean case of emergence through repeated coupling. The local events
are simple: same location, same cue, same entry action. The higher-level
result is a framework in which study is no longer reconstructed from scratch.
The evidence lane here is \textbf{medium}: complexity and emergence texts plus
the metastability review justify the structure of the explanation, but no
strong empirical paper in the current source set directly proves this exact
study ritual.
\end{discussion}

\subsection{Example 2: Phone Checking As Interference Architecture}

\begin{example}[Competing route gains historical priority]
Suppose a person repeatedly checks the phone at every short idle boundary:
after a paragraph, after a meal, while standing in line, before sleep. Over
time the idle-to-phone route becomes faster, more expected, and more
self-reinforcing than idle-to-work resumption.
\end{example}

\begin{discussion}
This example shows interference rather than simple habit accumulation. The old
route now intrudes into unrelated contexts because it has acquired superior
re-entry economy. In system terms, the phone route has become a highly
accessible neighbouring basin. The explanatory value is \textbf{medium} and
comes from path dependence, emergence, and attractor language rather than from
any claim of direct neural decoding.
\end{discussion}

\subsection{Example 3: Exercise Identity Through Multiscale Organization}

\begin{example}[Framework forms across body, schedule, and social timing]
A worker begins going to the gym on three fixed evenings, lays out clothes in
advance, travels with the same route, and meets the same partner at the same
time. After several weeks, the ``exercise evening'' no longer feels like a
fresh moral decision each time.
\end{example}

\begin{discussion}
This case is important because it is multiscale. The framework is not formed by
inner resolve alone. Body preparation, route repetition, calendar structure,
and social synchronization all contribute. The macro-pattern becomes
explanatory because it changes what the evening \emph{is}. The system now
carries a regime, not merely a wish.
\end{discussion}

\subsection{Example 4: Research Writing And Negative Path Dependence}

\begin{example}[Avoidance becomes organized competence]
A graduate student repeatedly resolves uncertainty by reading one more paper,
making one more note, or polishing citations rather than writing draft prose.
Eventually the pre-writing state becomes highly skilled, while visible writing
feels increasingly costly.
\end{example}

\begin{discussion}
This is a good example of negative emergence. The system does not merely fail
to write; it becomes organized around avoiding exposure. The resulting
framework has explanatory value because it changes later cost structure. The
manuscript should therefore not ask only ``Why is the student reluctant
today?'' but also ``What repeated micro-policy has trained this reluctance into
a stable route?''
\end{discussion}

\subsection{Example 5: Group Norm Formation}

\begin{example}[Frameworks can be collective]
In a team meeting, people repeatedly reward quick agreement and implicitly
punish slow clarification. After enough repetitions, the group develops a
framework in which fast consensus feels natural and careful dissent feels
socially expensive.
\end{example}

\begin{discussion}
This example broadens the book's relevance. Framework formation is not only
personal. Repeated local interactions can generate collective operating
patterns that later govern what individuals experience as easy or costly.
Complexity texts and emergence theory are especially valuable here, because the
macro-pattern cannot be reduced to any one participant's isolated intention.
\end{discussion}

\section{Evidence Boundaries}

This chapter must remain strict about what it has and has not established.

\begin{enumerate}[nosep]
  \item \textbf{Strong:} organized biological states, stability structure, and
        nontrivial transition dynamics are legitimate scientific topics, as
        established by the empirical consciousness chapter.
  \item \textbf{Medium:} metastability, causal emergence, attractor language,
        and complex-systems reasoning provide a defensible explanatory bridge
        for discussing framework formation.
  \item \textbf{Speculative if overextended:} any claim that a specific daily
        routine is generated by the exact same mechanism as a gene regulatory
        network, or that emergence language alone proves subjective experience
        or free will.
\end{enumerate}

\begin{remark}
This evidence ladder is not a weakness. It is what keeps the chapter from
pretending that sophisticated vocabulary is the same thing as strong proof.
\end{remark}

\section{Why Formation Matters For Intervention}

Once frameworks are understood as emergent historical structures, intervention
design changes. The central question is no longer only ``How do I resist the
bad state now?'' It becomes ``Which repeated local couplings are still
manufacturing the bad state, and which loop redesign will produce a different
macro-pattern later?''

This gives a more rigorous rationale for intervention:
\begin{enumerate}[nosep]
  \item identify the repeated cue-action-reward loop,
  \item identify which scale is carrying the regime,
  \item identify whether the regime is rigid, metastable, or fragile,
  \item redesign repeated conditions rather than relying only on heroic
        moment-level resistance.
\end{enumerate}

That is the business value of emergence language. It turns the manuscript from
a theory of resisting bad states into a theory of manufacturing better ones.

\section{Recommended Reading}

\begin{readingbox}
\textbf{Foundation}
\begin{itemize}[nosep]
  \item John Holland, \emph{Emergence}.
  \item \emph{Introduction to the Theory of Complex Systems}.
  \item \emph{Complexity Science}, for a broader systems vocabulary around
        self-organization and multiscale order.
\end{itemize}
\textbf{Biological and Review Layer}
\begin{itemize}[nosep]
  \item \emph{Metastability demystified: the foundational past, the pragmatic
        present and the promising future} (Nature Reviews Neuroscience, 2024,
        medium).
  \item \emph{Associative conditioning in gene regulatory network models
        increases integrative causal emergence} (Communications Biology, 2025,
        medium).
  \item \emph{Cortical connectivity, local dynamics and stability correlates of
        global conscious states} (Communications Biology, 2025, strong as
        biological background, medium for this chapter's transfer claims).
\end{itemize}
\textbf{Mathematical bridge}
\begin{itemize}[nosep]
  \item Eugene Izhikevich, \emph{Dynamical Systems in Neuroscience}.
  \item Review literature on attractor and integrator networks in the brain
        (Nature Reviews Neuroscience, 2022, medium).
\end{itemize}
\textbf{Reading discipline}
\begin{itemize}[nosep]
  \item Use the biology and review layer to justify organized state formation
        and stability.
  \item Use the systems and mathematics layer to clarify what kind of higher-
        level pattern counts as emergent.
  \item Do not treat frontier consciousness-physics proposals as evidence for
        this chapter's core claims.
\end{itemize}
\end{readingbox}

\begin{summarybox}
This chapter now gives a stricter account of where frameworks come from. Its
source-backed claim is not that emergence has been fully measured for every
daily routine, but that repeated local interactions, conditioning history,
multiscale coupling, and metastable organization can generate higher-level
regimes that later restructure action costs. Biology gives the chapter
legitimacy, systems language gives it conceptual reach, and mathematical
structure prevents it from collapsing into metaphor. That is what upgrades the
manuscript from a static account of persistence into a developmental account of
formation, consolidation, interference, and redesign.
\end{summarybox}
